We wish to express our deepest gratitude to our principal supervisor, Prof.
Harinda Fernando, and our co-supervisor, Mr. Indunil Daluwatte, for their
invaluable guidance, steady support and constructive feedback throughout this
research. Their expertise and encouragement shaped the direction, development
and quality of this work, and each individual component of the NeuroShell
framework.

We are also sincerely grateful to the academic staff of the Faculty of Computing
at the Sri Lanka Institute of Information Technology (SLIIT) for the knowledge,
the supportive learning environment and the resources that made research in
cybersecurity possible. The questionnaires, evaluation instruments and lab
facilities provided by the department supported every phase of this project.

Our special thanks go to the members of Research Group R26-CS-018 for their
cooperation, their technical discussions and the spirit of collaboration in
which we built the NeuroShell framework together. Integrating four individually
developed components into a single orchestration chain demanded sustained
communication about interface contracts, data formats and shared state, and this
work would not have been possible without that collective effort.

Finally, we thank our families and friends from the bottom of our hearts for
their patience, encouragement and moral support throughout our B.Sc. (Hons)
degree programme and the completion of this research project.